% TOP_PROBLEM: {"id": 30007151, "title": "Nonnegative Einstein curvature classification in dimension four", "category_id": 6, "category": {"id": 6, "name": "geometry", "display_name": "Geometry", "description": "Euclidean and non-Euclidean geometry, geometric structures.", "slug": "geometry", "order_index": 6, "created_at": "2026-07-31T15:26:25.671Z"}, "status": "open"}
% FIRST_POSTED: 2026-10-06
% ATTEMPT_STATUS: solved
% OPENAI_CLAIM: {"schema_version":1,"source_repo":"https://github.com/openai/math","source_commit":"adc7f1241b42e322a6451854ab7e4b4c146bf78a","release_date":"2026-10-06","match":"full","resolution":"proved","independently_reviewed":false,"second_pass":"agreed","families":[{"family":"348","title":"Nonnegative-curvature Einstein classification and an <i>L</i><sup>2</sup> topological gap","subject":"Differential geometry","manuscripts":[{"title":"Zero-Plane Rigidity for Einstein Four-Manifolds","pdf_path":"preprints/Zero-Plane-Rigidity-for-Einstein-Four-Manifolds-October-4-2026/einstein-boundary.pdf","readme_path":"preprints/Zero-Plane-Rigidity-for-Einstein-Four-Manifolds-October-4-2026/README.md","date":"2026-10-04","theorem_ref":"preprints/Zero-Plane-Rigidity-for-Einstein-Four-Manifolds-October-4-2026/build/sections/01-introduction.tex#cor:nonnegative-classification"}],"lean":{"doc_path":"lean/docs/348.md","comparators":[],"declarations":[],"covers_main_theorem":false},"reasoning_trace":null}]}
\documentclass[11pt]{article}
\usepackage[margin=25mm]{geometry}
\usepackage{fontspec}
\setmainfont{Latin Modern Roman}
\usepackage{amsmath,amssymb,mathtools}
\usepackage{unicode-math}
\setmathfont{Latin Modern Math}
\usepackage{enumitem}
\usepackage{xurl,hyperref,bookmark}
\hypersetup{colorlinks=true,urlcolor=blue}
\setcounter{secnumdepth}{0}
\setlength{\parindent}{0pt}
\setlength{\parskip}{5pt}
\setlength{\emergencystretch}{3em}
\begin{document}
\section*{Nonnegative Einstein curvature classification in dimension four}

\subsection{OpenAI's claim}
\textit{Zero-plane Einstein rigidity}: The claim assumes a positive Einstein constant on a closed four-manifold. If sectional curvature is nonnegative and one plane is flat, the universal cover is the product of two round constant-curvature surfaces. The radii are fixed by the Einstein normalization.

\subsection{Outline of the argument}
\textit{Zero-plane Einstein rigidity}: Bonnet--Myers reduces the analysis to the compact simply connected cover. The two Weyl blocks satisfy a sharp coupled spectral constraint at the boundary of the curvature cone. Extending the weighted comparisons to this boundary and tracking equality forces the product geometry rather than merely a weaker norm bound.

\subsection{Scope relative to this problem}
\textit{Zero-plane Einstein rigidity}: The claim assumes a positive Einstein constant on a closed four-manifold. If sectional curvature is nonnegative and one plane is flat, the universal cover is the product of two round constant-curvature surfaces. The radii are fixed by the Einstein normalization. The parameters, regularity assumptions and normalization in this question are retained in applying that assertion.

With the hypotheses and conventions above, the claimed result addresses the whole mathematical assertion on this page.

\subsection{Formal verification}
Family 348 describes Lean coverage for the strict-positive-curvature Einstein classification; zero-plane rigidity and the small-energy gap are separate.

\textit{Zero-plane Einstein rigidity}: No manuscript-specific Comparator is listed in the pinned inventory.

The documented selection does not cover the entire manuscript assertion relevant here.

This site has not run the Lean development. These links describe OpenAI’s selected formal statements and should be read with their stated scope.

\subsection{Status caveat}
OpenAI reports that an internal model produced the manuscripts. Their arguments have not been independently verified as proofs by this site. Any portion outside the documented formal scope may contain mistakes, and the catalogue label reports the announced mathematical claim.

\subsection{Sources}
\begin{itemize}
\item OpenAI, \emph{Zero-Plane Rigidity for Einstein Four-Manifolds}, 2026-10-04; citation key \texttt{OAI:Zero-Plane-Rigidity-for-Einstein-Four-Manifolds-October-4-2026}. \href{https://github.com/openai/math/blob/adc7f1241b42e322a6451854ab7e4b4c146bf78a/preprints/Zero-Plane-Rigidity-for-Einstein-Four-Manifolds-October-4-2026/einstein-boundary.pdf}{PDF (GitHub)}; \href{https://github.com/openai/math/raw/adc7f1241b42e322a6451854ab7e4b4c146bf78a/preprints/Zero-Plane-Rigidity-for-Einstein-Four-Manifolds-October-4-2026/einstein-boundary.pdf}{PDF (direct)}; \href{https://github.com/openai/math/blob/adc7f1241b42e322a6451854ab7e4b4c146bf78a/preprints/Zero-Plane-Rigidity-for-Einstein-Four-Manifolds-October-4-2026/README.md}{Manuscript page and citation}; \href{https://github.com/openai/math/blob/main/preprints/Zero-Plane-Rigidity-for-Einstein-Four-Manifolds-October-4-2026/einstein-boundary.pdf}{Latest version}.
\item Family 348: \href{https://github.com/openai/math/blob/adc7f1241b42e322a6451854ab7e4b4c146bf78a/lean/docs/348.md}{Lean scope}.
\item \href{https://github.com/openai/math/blob/adc7f1241b42e322a6451854ab7e4b4c146bf78a/CONTENTS.md}{OpenAI catalogue entry}.
\end{itemize}
\end{document}
